\documentclass[11pt]{article}

\usepackage[margin=1in]{geometry}
\usepackage{amsmath,amssymb}
\usepackage{booktabs}
\usepackage{array}
\usepackage{hyperref}

\title{From Barwick et al. (2024) to the TidalLanes Congestion Parameter}
\author{}
\date{\today}

\begin{document}
\maketitle

\section{Objective}

This note shows how to take the speed--density elasticity estimated by
Barwick et al.\ (AER 2024) on Beijing traffic data and convert it into the
congestion parameter that TidalLanes actually uses. The conversion has two
steps, corresponding to the two structural elasticities in the model:
\begin{itemize}
\item $\delta_1$: the \emph{congestion elasticity}. It governs how traffic
density raises physical travel time (equivalently, lowers speed). This is the
object Barwick et al.\ estimate.
\item $\delta_0=1/\theta$: the \emph{cost-curvature conversion}. It governs how
physical travel time is turned into an iceberg trade cost, and is pinned down
by the Frechet shape parameter $\theta$, not by any traffic data.
\end{itemize}
The parameter that finally enters the iceberg transportation-cost equation is
the \emph{product} of the two:
\begin{equation}
\lambda_{\mathrm{AA}}=\delta_0\,\delta_1=\frac{\delta_1}{\theta}.
\end{equation}
So the empirical paper only supplies $\delta_1$; the factor $1/\theta$ comes
from the model's own calibration.

\section{Notation}

Let $S_{kl}$ be speed and $D_{kl}$ traffic density on link $(k,l)$. The Beijing
papers define density as vehicles per lane-kilometre,
\begin{equation}
D_{kl}=\frac{\xi_{kl}}{n_{kl}\,\ell_{kl}},
\end{equation}
where $\xi_{kl}$ is directed link traffic, $n_{kl}$ is the directed lane count,
and $\ell_{kl}$ is link length. Physical travel time on the link is
\begin{equation}
\tau_{kl}=\frac{\ell_{kl}}{S_{kl}}.
\end{equation}
Holding link length fixed, the elasticity of travel time with respect to
density equals minus the elasticity of speed with respect to density.

\section{The Barwick et al.\ estimate}

Barwick et al.\ estimate a log-log speed--density relationship,
\begin{equation}
\ln S_{st}=e\,\ln D_{st}+X_{st}'\beta+\varepsilon_{st},
\end{equation}
with a city-wide IV estimate
\begin{equation}
\hat{e}\approx -1.1.
\end{equation}
Because the equation is already in logs, $\hat{e}$ \emph{is} the speed--density
elasticity directly --- no baseline $(D,S)$ point is needed. (This is the
practical advantage over Yang et al.\ (2020), whose coefficient is a level
slope and would first have to be evaluated at a chosen baseline.)

\section{Step 1: congestion elasticity $\delta_1=\lambda_{\mathrm{phys}}$}

The congestion elasticity is defined on physical travel time, or equivalently
on inverse speed $S^{-1}$. Since $\ln S^{-1}=-\ln S$,
\begin{equation}
\delta_1
=\frac{\partial \ln S^{-1}}{\partial \ln D}
=-\frac{\partial \ln S}{\partial \ln D}
=-\hat{e}\approx 1.1.
\end{equation}
In the physical-time version of the TidalLanes congestion technology,
\begin{equation}
\tau_{kl}=\bar{\tau}_{kl}\,D_{kl}^{\lambda_{\mathrm{phys}}},
\qquad
\lambda_{\mathrm{phys}}=-\frac{\partial \ln S_{kl}}{\partial \ln D_{kl}}
=-\hat{e}\approx 1.1 .
\end{equation}
So $\lambda_{\mathrm{phys}}\equiv\delta_1\approx 1.1$.

\section{Step 2: cost-curvature conversion $\delta_0=1/\theta$}

If TidalLanes follows the Allen--Arkolakis iceberg-cost convention, the trade
cost is a power of physical travel time,
\begin{equation}
t_{kl}=\tau_{kl}^{1/\theta}.
\end{equation}
Substituting the congestion technology gives
\begin{equation}
t_{kl}
=\left(\bar{\tau}_{kl}\,D_{kl}^{\lambda_{\mathrm{phys}}}\right)^{1/\theta}
=\bar{t}_{kl}\,D_{kl}^{\lambda_{\mathrm{AA}}},
\qquad
\lambda_{\mathrm{AA}}=\frac{\lambda_{\mathrm{phys}}}{\theta}=\frac{\delta_1}{\theta}.
\end{equation}
With the current calibration $\theta=6.83$,
\begin{equation}
\lambda_{\mathrm{AA}}=\frac{1.1}{6.83}\approx 0.161 .
\end{equation}

\section{Summary}

\renewcommand{\arraystretch}{1.3}
\begin{center}
\begin{tabular}{@{}>{\raggedright\arraybackslash}p{0.30\textwidth}
                  >{\raggedright\arraybackslash}p{0.32\textwidth}
                  >{\raggedright\arraybackslash}p{0.28\textwidth}@{}}
\toprule
Quantity & Expression & Value \\
\midrule
Barwick speed--density elasticity & $\hat{e}$ & $-1.1$ \\
Congestion elasticity (physical time) & $\delta_1=\lambda_{\mathrm{phys}}=-\hat{e}$ & $1.1$ \\
Cost-curvature conversion & $\delta_0=1/\theta$ & $1/6.83\approx 0.146$ \\
Model congestion parameter (iceberg) & $\lambda_{\mathrm{AA}}=\delta_1/\theta$ & $1.1/6.83\approx 0.161$ \\
\bottomrule
\end{tabular}
\end{center}

\noindent Which number to plug in depends on how the model is written:
\begin{itemize}
\item physical-travel-time form $\tau=\bar{\tau}D^{\lambda_{\mathrm{phys}}}$:
use $\lambda_{\mathrm{phys}}\approx 1.1$;
\item iceberg-cost form $t=\bar{t}D^{\lambda_{\mathrm{AA}}}$:
use $\lambda_{\mathrm{AA}}\approx 0.161$ (at $\theta=6.83$).
\end{itemize}
For reference, the current code uses $\lambda_{\mathrm{code}}=0.488/6.83\approx
0.071$, which corresponds to a congestion elasticity of $\delta_1=0.488$ from
the AA Seattle calibration --- not the Beijing-specific Barwick value.

\section{One caveat on the density regressor}

The Barwick elasticity is measured against density in vehicles per
\emph{lane-km}, $D_{kl}=\xi_{kl}/(n_{kl}\ell_{kl})$. As long as the TidalLanes
congestion regressor is also per-lane-km (or per-lane with link length fixed
within each link, in which case $\ln\ell_{kl}$ is absorbed into the constant
$\bar{\tau}_{kl}$), the elasticity carries over unchanged. If the regressor
omits link length in a way that varies across links, the constant term shifts
but the elasticity $\delta_1$ is unaffected.

\end{document}
